% =====================================================================
%  phuluc_nguon.tex  —  PHỤ LỤC: BẢNG ĐỐI CHIẾU LUẬN ĐIỂM - NGUỒN TRÍCH DẪN
%  Gói cần thêm vào preamble:  \usepackage{longtable}
%  Chèn vào main.tex sau \appendix (hoặc sau chương 5) bằng:
%      % =====================================================================
%  phuluc_nguon.tex  —  PHỤ LỤC: BẢNG ĐỐI CHIẾU LUẬN ĐIỂM - NGUỒN TRÍCH DẪN
%  Gói cần thêm vào preamble:  \usepackage{longtable}
%  Chèn vào main.tex sau \appendix (hoặc sau chương 5) bằng:
%      % =====================================================================
%  phuluc_nguon.tex  —  PHỤ LỤC: BẢNG ĐỐI CHIẾU LUẬN ĐIỂM - NGUỒN TRÍCH DẪN
%  Gói cần thêm vào preamble:  \usepackage{longtable}
%  Chèn vào main.tex sau \appendix (hoặc sau chương 5) bằng:
%      % =====================================================================
%  phuluc_nguon.tex  —  PHỤ LỤC: BẢNG ĐỐI CHIẾU LUẬN ĐIỂM - NGUỒN TRÍCH DẪN
%  Gói cần thêm vào preamble:  \usepackage{longtable}
%  Chèn vào main.tex sau \appendix (hoặc sau chương 5) bằng:
%      \input{phuluc_nguon}
% =====================================================================

\chapter*{PHỤ LỤC A. ĐỐI CHIẾU LUẬN ĐIỂM VÀ NGUỒN TRÍCH DẪN}
\addcontentsline{toc}{chapter}{PHỤ LỤC A. ĐỐI CHIẾU LUẬN ĐIỂM VÀ NGUỒN TRÍCH DẪN}
\label{app:nguon}

Phụ lục này liệt kê từng luận điểm lý thuyết được sử dụng trong khóa luận kèm
nguồn gốc học thuật tương ứng, nhằm bảo đảm mọi khẳng định đều truy nguyên được
về tài liệu công bố chính thống.

\begin{longtable}{|p{0.06\textwidth}|p{0.44\textwidth}|p{0.38\textwidth}|}
\hline
\textbf{Mục} & \textbf{Luận điểm} & \textbf{Nguồn} \\
\hline
\endfirsthead
\hline
\textbf{Mục} & \textbf{Luận điểm} & \textbf{Nguồn} \\
\hline
\endhead

\multicolumn{3}{l}{\textit{A.1. Lý thuyết tồn kho}} \\
\hline
\ref{subsec:thuctrang} & Quản lý tồn kho đa điểm, đa mặt hàng trong chuỗi cung ứng bán lẻ &
\textcite{silver2016}; \textcite{zipkin2000} \\

\ref{subsec:thuctrang} & Mô tả bộ dữ liệu M5: quy mô, cấu trúc phân cấp, giai đoạn thu thập &
\textcite{makridakis2022background}; \textcite{m5dataset} \\

\ref{subsec:thuctrang} & Tiêu chí phân loại dạng nhu cầu (ngưỡng $p = 1{,}32$; $CV^2 = 0{,}49$) &
\textcite{syntetos2005categorization} \\

\ref{subsec:thuctrang} & Nhu cầu gián đoạn và hạn chế của phương pháp dự báo thông thường &
\textcite{croston1972}; \textcite{syntetos2005accuracy} \\

\ref{subsec:thachthuc} & Bất định của nhu cầu và thời gian giao hàng; hiện tượng vượt đơn &
\textcite{zipkin2000}; \textcite{silver2016} \\

\ref{subsec:thachthuc} & Cấu trúc chi phí tồn kho và bài toán đánh đổi &
\textcite{zipkin2000} \\

\ref{subsec:chinhSachCoDien} & Mô hình lượng đặt hàng kinh tế (EOQ) &
\textcite{harris1913} \\

\ref{subsec:chinhSachCoDien} & Tính tối ưu của chính sách $(s, S)$ và các điều kiện kèm theo &
\textcite{scarf1960} \\

\ref{subsec:chinhSachCoDien} & Mô hình người bán báo (Newsvendor) một chu kỳ &
\textcite{arrow1951} \\

\ref{subsec:moHinhTonKho} & Bài toán mất đơn hàng không có nghiệm giải tích dạng đóng &
\textcite{zipkin2008} \\

\ref{subsec:phamvingoai} & Cấu trúc tồn kho đa cấp (nằm ngoài phạm vi) &
\textcite{clark1960} \\

\ref{subsec:phamvingoai} & Hiệu ứng bullwhip (nằm ngoài phạm vi) &
\textcite{lee1997} \\
\hline

\multicolumn{3}{l}{\textit{A.2. Học tăng cường}} \\
\hline
\ref{subsec:mdp} & Quá trình quyết định Markov và quy hoạch động ngẫu nhiên &
\textcite{puterman1994} \\

\ref{subsec:hamGiaTri} & Hàm giá trị, phương trình Bellman, hàm lợi thế, định lý policy gradient &
\textcite{sutton2018} \\

\ref{subsec:gae} & Ước lượng lợi thế tổng quát (GAE) và đánh đổi bias--variance &
\textcite{schulman2016gae} \\

\ref{subsec:ppo} & Vùng tin cậy trong tối ưu hóa chính sách &
\textcite{schulman2015trpo} \\

\ref{subsec:ppo} & Hàm mục tiêu cắt của PPO &
\textcite{schulman2017ppo} \\

\ref{subsec:rewardDecomposition} & Điều kiện bảo toàn chính sách tối ưu khi định hình phần thưởng &
\textcite{ng1999} \\
\hline

\multicolumn{3}{l}{\textit{A.3. Học tăng cường đa tác tử}} \\
\hline
\ref{subsec:decPomdp} & Độ phức tạp NEXP-đầy đủ của điều khiển phi tập trung &
\textcite{bernstein2002} \\

\ref{subsec:decPomdp} & Mô hình Dec-POMDP &
\textcite{oliehoek2016} \\

\ref{subsec:hocDocLap} & Học độc lập và vấn đề môi trường không dừng &
\textcite{tan1993} \\

\ref{subsec:ippo} & Thuật toán IPPO &
\textcite{dewitt2020ippo} \\

\ref{subsec:ippo} & MAPPO và khung huấn luyện tập trung, thực thi phân tán &
\textcite{yu2022mappo} \\

\ref{subsec:paramSharing} & Chia sẻ tham số giữa các tác tử đồng nhất &
\textcite{gupta2017} \\

\ref{subsec:rewardDecomposition} & Bài toán gán tín dụng đa tác tử; phần thưởng chênh lệch (so sánh với phần thưởng phân rã cục bộ) &
\textcite{wolpert2001} \\
\hline

\multicolumn{3}{l}{\textit{A.4. Ứng dụng, môi trường và giao thức đánh giá}} \\
\hline
\ref{subsec:rlTonKho} & Học tăng cường sâu cho trò chơi phân phối bia &
\textcite{oroojlooyjadid2022} \\

\ref{subsec:rlTonKho} & Đánh giá học tăng cường sâu trên các lớp bài toán tồn kho &
\textcite{gijsbrechts2022} \\

\ref{sec:moiTruong} & Chuẩn giao diện môi trường học tăng cường &
\textcite{brockman2016gym}; \textcite{towers2024gymnasium} \\

\ref{subsec:daHatGiong} & Yêu cầu đa hạt giống và tính tái lập trong đánh giá &
\textcite{henderson2018} \\

\ref{subsec:daHatGiong} & Phương pháp báo cáo kết quả có kiểm định thống kê &
\textcite{agarwal2021} \\
\hline
\end{longtable}

\vspace{0.5\baselineskip}
\noindent\textbf{Ghi chú.} Các số liệu thống kê mô tả bộ dữ liệu M5 trình bày
trong Mục~\ref{subsec:thuctrang} và Bảng~\ref{tab:phanloainhucau} do nhóm tác
giả tự tính toán từ hai tệp \texttt{sales\_train\_evaluation.csv} và
\texttt{calendar.csv}; mã nguồn tính toán được đính kèm trong kho mã của khóa
luận.

% =====================================================================

\chapter*{PHỤ LỤC A. ĐỐI CHIẾU LUẬN ĐIỂM VÀ NGUỒN TRÍCH DẪN}
\addcontentsline{toc}{chapter}{PHỤ LỤC A. ĐỐI CHIẾU LUẬN ĐIỂM VÀ NGUỒN TRÍCH DẪN}
\label{app:nguon}

Phụ lục này liệt kê từng luận điểm lý thuyết được sử dụng trong khóa luận kèm
nguồn gốc học thuật tương ứng, nhằm bảo đảm mọi khẳng định đều truy nguyên được
về tài liệu công bố chính thống.

\begin{longtable}{|p{0.06\textwidth}|p{0.44\textwidth}|p{0.38\textwidth}|}
\hline
\textbf{Mục} & \textbf{Luận điểm} & \textbf{Nguồn} \\
\hline
\endfirsthead
\hline
\textbf{Mục} & \textbf{Luận điểm} & \textbf{Nguồn} \\
\hline
\endhead

\multicolumn{3}{l}{\textit{A.1. Lý thuyết tồn kho}} \\
\hline
\ref{subsec:thuctrang} & Quản lý tồn kho đa điểm, đa mặt hàng trong chuỗi cung ứng bán lẻ &
\textcite{silver2016}; \textcite{zipkin2000} \\

\ref{subsec:thuctrang} & Mô tả bộ dữ liệu M5: quy mô, cấu trúc phân cấp, giai đoạn thu thập &
\textcite{makridakis2022background}; \textcite{m5dataset} \\

\ref{subsec:thuctrang} & Tiêu chí phân loại dạng nhu cầu (ngưỡng $p = 1{,}32$; $CV^2 = 0{,}49$) &
\textcite{syntetos2005categorization} \\

\ref{subsec:thuctrang} & Nhu cầu gián đoạn và hạn chế của phương pháp dự báo thông thường &
\textcite{croston1972}; \textcite{syntetos2005accuracy} \\

\ref{subsec:thachthuc} & Bất định của nhu cầu và thời gian giao hàng; hiện tượng vượt đơn &
\textcite{zipkin2000}; \textcite{silver2016} \\

\ref{subsec:thachthuc} & Cấu trúc chi phí tồn kho và bài toán đánh đổi &
\textcite{zipkin2000} \\

\ref{subsec:chinhSachCoDien} & Mô hình lượng đặt hàng kinh tế (EOQ) &
\textcite{harris1913} \\

\ref{subsec:chinhSachCoDien} & Tính tối ưu của chính sách $(s, S)$ và các điều kiện kèm theo &
\textcite{scarf1960} \\

\ref{subsec:chinhSachCoDien} & Mô hình người bán báo (Newsvendor) một chu kỳ &
\textcite{arrow1951} \\

\ref{subsec:moHinhTonKho} & Bài toán mất đơn hàng không có nghiệm giải tích dạng đóng &
\textcite{zipkin2008} \\

\ref{subsec:phamvingoai} & Cấu trúc tồn kho đa cấp (nằm ngoài phạm vi) &
\textcite{clark1960} \\

\ref{subsec:phamvingoai} & Hiệu ứng bullwhip (nằm ngoài phạm vi) &
\textcite{lee1997} \\
\hline

\multicolumn{3}{l}{\textit{A.2. Học tăng cường}} \\
\hline
\ref{subsec:mdp} & Quá trình quyết định Markov và quy hoạch động ngẫu nhiên &
\textcite{puterman1994} \\

\ref{subsec:hamGiaTri} & Hàm giá trị, phương trình Bellman, hàm lợi thế, định lý policy gradient &
\textcite{sutton2018} \\

\ref{subsec:gae} & Ước lượng lợi thế tổng quát (GAE) và đánh đổi bias--variance &
\textcite{schulman2016gae} \\

\ref{subsec:ppo} & Vùng tin cậy trong tối ưu hóa chính sách &
\textcite{schulman2015trpo} \\

\ref{subsec:ppo} & Hàm mục tiêu cắt của PPO &
\textcite{schulman2017ppo} \\

\ref{subsec:rewardDecomposition} & Điều kiện bảo toàn chính sách tối ưu khi định hình phần thưởng &
\textcite{ng1999} \\
\hline

\multicolumn{3}{l}{\textit{A.3. Học tăng cường đa tác tử}} \\
\hline
\ref{subsec:decPomdp} & Độ phức tạp NEXP-đầy đủ của điều khiển phi tập trung &
\textcite{bernstein2002} \\

\ref{subsec:decPomdp} & Mô hình Dec-POMDP &
\textcite{oliehoek2016} \\

\ref{subsec:hocDocLap} & Học độc lập và vấn đề môi trường không dừng &
\textcite{tan1993} \\

\ref{subsec:ippo} & Thuật toán IPPO &
\textcite{dewitt2020ippo} \\

\ref{subsec:ippo} & MAPPO và khung huấn luyện tập trung, thực thi phân tán &
\textcite{yu2022mappo} \\

\ref{subsec:paramSharing} & Chia sẻ tham số giữa các tác tử đồng nhất &
\textcite{gupta2017} \\

\ref{subsec:rewardDecomposition} & Bài toán gán tín dụng đa tác tử; phần thưởng chênh lệch (so sánh với phần thưởng phân rã cục bộ) &
\textcite{wolpert2001} \\
\hline

\multicolumn{3}{l}{\textit{A.4. Ứng dụng, môi trường và giao thức đánh giá}} \\
\hline
\ref{subsec:rlTonKho} & Học tăng cường sâu cho trò chơi phân phối bia &
\textcite{oroojlooyjadid2022} \\

\ref{subsec:rlTonKho} & Đánh giá học tăng cường sâu trên các lớp bài toán tồn kho &
\textcite{gijsbrechts2022} \\

\ref{sec:moiTruong} & Chuẩn giao diện môi trường học tăng cường &
\textcite{brockman2016gym}; \textcite{towers2024gymnasium} \\

\ref{subsec:daHatGiong} & Yêu cầu đa hạt giống và tính tái lập trong đánh giá &
\textcite{henderson2018} \\

\ref{subsec:daHatGiong} & Phương pháp báo cáo kết quả có kiểm định thống kê &
\textcite{agarwal2021} \\
\hline
\end{longtable}

\vspace{0.5\baselineskip}
\noindent\textbf{Ghi chú.} Các số liệu thống kê mô tả bộ dữ liệu M5 trình bày
trong Mục~\ref{subsec:thuctrang} và Bảng~\ref{tab:phanloainhucau} do nhóm tác
giả tự tính toán từ hai tệp \texttt{sales\_train\_evaluation.csv} và
\texttt{calendar.csv}; mã nguồn tính toán được đính kèm trong kho mã của khóa
luận.

% =====================================================================

\chapter*{PHỤ LỤC A. ĐỐI CHIẾU LUẬN ĐIỂM VÀ NGUỒN TRÍCH DẪN}
\addcontentsline{toc}{chapter}{PHỤ LỤC A. ĐỐI CHIẾU LUẬN ĐIỂM VÀ NGUỒN TRÍCH DẪN}
\label{app:nguon}

Phụ lục này liệt kê từng luận điểm lý thuyết được sử dụng trong khóa luận kèm
nguồn gốc học thuật tương ứng, nhằm bảo đảm mọi khẳng định đều truy nguyên được
về tài liệu công bố chính thống.

\begin{longtable}{|p{0.06\textwidth}|p{0.44\textwidth}|p{0.38\textwidth}|}
\hline
\textbf{Mục} & \textbf{Luận điểm} & \textbf{Nguồn} \\
\hline
\endfirsthead
\hline
\textbf{Mục} & \textbf{Luận điểm} & \textbf{Nguồn} \\
\hline
\endhead

\multicolumn{3}{l}{\textit{A.1. Lý thuyết tồn kho}} \\
\hline
\ref{subsec:thuctrang} & Quản lý tồn kho đa điểm, đa mặt hàng trong chuỗi cung ứng bán lẻ &
\textcite{silver2016}; \textcite{zipkin2000} \\

\ref{subsec:thuctrang} & Mô tả bộ dữ liệu M5: quy mô, cấu trúc phân cấp, giai đoạn thu thập &
\textcite{makridakis2022background}; \textcite{m5dataset} \\

\ref{subsec:thuctrang} & Tiêu chí phân loại dạng nhu cầu (ngưỡng $p = 1{,}32$; $CV^2 = 0{,}49$) &
\textcite{syntetos2005categorization} \\

\ref{subsec:thuctrang} & Nhu cầu gián đoạn và hạn chế của phương pháp dự báo thông thường &
\textcite{croston1972}; \textcite{syntetos2005accuracy} \\

\ref{subsec:thachthuc} & Bất định của nhu cầu và thời gian giao hàng; hiện tượng vượt đơn &
\textcite{zipkin2000}; \textcite{silver2016} \\

\ref{subsec:thachthuc} & Cấu trúc chi phí tồn kho và bài toán đánh đổi &
\textcite{zipkin2000} \\

\ref{subsec:chinhSachCoDien} & Mô hình lượng đặt hàng kinh tế (EOQ) &
\textcite{harris1913} \\

\ref{subsec:chinhSachCoDien} & Tính tối ưu của chính sách $(s, S)$ và các điều kiện kèm theo &
\textcite{scarf1960} \\

\ref{subsec:chinhSachCoDien} & Mô hình người bán báo (Newsvendor) một chu kỳ &
\textcite{arrow1951} \\

\ref{subsec:moHinhTonKho} & Bài toán mất đơn hàng không có nghiệm giải tích dạng đóng &
\textcite{zipkin2008} \\

\ref{subsec:phamvingoai} & Cấu trúc tồn kho đa cấp (nằm ngoài phạm vi) &
\textcite{clark1960} \\

\ref{subsec:phamvingoai} & Hiệu ứng bullwhip (nằm ngoài phạm vi) &
\textcite{lee1997} \\
\hline

\multicolumn{3}{l}{\textit{A.2. Học tăng cường}} \\
\hline
\ref{subsec:mdp} & Quá trình quyết định Markov và quy hoạch động ngẫu nhiên &
\textcite{puterman1994} \\

\ref{subsec:hamGiaTri} & Hàm giá trị, phương trình Bellman, hàm lợi thế, định lý policy gradient &
\textcite{sutton2018} \\

\ref{subsec:gae} & Ước lượng lợi thế tổng quát (GAE) và đánh đổi bias--variance &
\textcite{schulman2016gae} \\

\ref{subsec:ppo} & Vùng tin cậy trong tối ưu hóa chính sách &
\textcite{schulman2015trpo} \\

\ref{subsec:ppo} & Hàm mục tiêu cắt của PPO &
\textcite{schulman2017ppo} \\

\ref{subsec:rewardDecomposition} & Điều kiện bảo toàn chính sách tối ưu khi định hình phần thưởng &
\textcite{ng1999} \\
\hline

\multicolumn{3}{l}{\textit{A.3. Học tăng cường đa tác tử}} \\
\hline
\ref{subsec:decPomdp} & Độ phức tạp NEXP-đầy đủ của điều khiển phi tập trung &
\textcite{bernstein2002} \\

\ref{subsec:decPomdp} & Mô hình Dec-POMDP &
\textcite{oliehoek2016} \\

\ref{subsec:hocDocLap} & Học độc lập và vấn đề môi trường không dừng &
\textcite{tan1993} \\

\ref{subsec:ippo} & Thuật toán IPPO &
\textcite{dewitt2020ippo} \\

\ref{subsec:ippo} & MAPPO và khung huấn luyện tập trung, thực thi phân tán &
\textcite{yu2022mappo} \\

\ref{subsec:paramSharing} & Chia sẻ tham số giữa các tác tử đồng nhất &
\textcite{gupta2017} \\

\ref{subsec:rewardDecomposition} & Bài toán gán tín dụng đa tác tử; phần thưởng chênh lệch (so sánh với phần thưởng phân rã cục bộ) &
\textcite{wolpert2001} \\
\hline

\multicolumn{3}{l}{\textit{A.4. Ứng dụng, môi trường và giao thức đánh giá}} \\
\hline
\ref{subsec:rlTonKho} & Học tăng cường sâu cho trò chơi phân phối bia &
\textcite{oroojlooyjadid2022} \\

\ref{subsec:rlTonKho} & Đánh giá học tăng cường sâu trên các lớp bài toán tồn kho &
\textcite{gijsbrechts2022} \\

\ref{sec:moiTruong} & Chuẩn giao diện môi trường học tăng cường &
\textcite{brockman2016gym}; \textcite{towers2024gymnasium} \\

\ref{subsec:daHatGiong} & Yêu cầu đa hạt giống và tính tái lập trong đánh giá &
\textcite{henderson2018} \\

\ref{subsec:daHatGiong} & Phương pháp báo cáo kết quả có kiểm định thống kê &
\textcite{agarwal2021} \\
\hline
\end{longtable}

\vspace{0.5\baselineskip}
\noindent\textbf{Ghi chú.} Các số liệu thống kê mô tả bộ dữ liệu M5 trình bày
trong Mục~\ref{subsec:thuctrang} và Bảng~\ref{tab:phanloainhucau} do nhóm tác
giả tự tính toán từ hai tệp \texttt{sales\_train\_evaluation.csv} và
\texttt{calendar.csv}; mã nguồn tính toán được đính kèm trong kho mã của khóa
luận.

% =====================================================================

\chapter*{PHỤ LỤC A. ĐỐI CHIẾU LUẬN ĐIỂM VÀ NGUỒN TRÍCH DẪN}
\addcontentsline{toc}{chapter}{PHỤ LỤC A. ĐỐI CHIẾU LUẬN ĐIỂM VÀ NGUỒN TRÍCH DẪN}
\label{app:nguon}

Phụ lục này liệt kê từng luận điểm lý thuyết được sử dụng trong khóa luận kèm
nguồn gốc học thuật tương ứng, nhằm bảo đảm mọi khẳng định đều truy nguyên được
về tài liệu công bố chính thống.

\begin{longtable}{|p{0.06\textwidth}|p{0.44\textwidth}|p{0.38\textwidth}|}
\hline
\textbf{Mục} & \textbf{Luận điểm} & \textbf{Nguồn} \\
\hline
\endfirsthead
\hline
\textbf{Mục} & \textbf{Luận điểm} & \textbf{Nguồn} \\
\hline
\endhead

\multicolumn{3}{l}{\textit{A.1. Lý thuyết tồn kho}} \\
\hline
\ref{subsec:thuctrang} & Quản lý tồn kho đa điểm, đa mặt hàng trong chuỗi cung ứng bán lẻ &
\textcite{silver2016}; \textcite{zipkin2000} \\

\ref{subsec:thuctrang} & Mô tả bộ dữ liệu M5: quy mô, cấu trúc phân cấp, giai đoạn thu thập &
\textcite{makridakis2022background}; \textcite{m5dataset} \\

\ref{subsec:thuctrang} & Tiêu chí phân loại dạng nhu cầu (ngưỡng $p = 1{,}32$; $CV^2 = 0{,}49$) &
\textcite{syntetos2005categorization} \\

\ref{subsec:thuctrang} & Nhu cầu gián đoạn và hạn chế của phương pháp dự báo thông thường &
\textcite{croston1972}; \textcite{syntetos2005accuracy} \\

\ref{subsec:thachthuc} & Bất định của nhu cầu và thời gian giao hàng; hiện tượng vượt đơn &
\textcite{zipkin2000}; \textcite{silver2016} \\

\ref{subsec:thachthuc} & Cấu trúc chi phí tồn kho và bài toán đánh đổi &
\textcite{zipkin2000} \\

\ref{subsec:chinhSachCoDien} & Mô hình lượng đặt hàng kinh tế (EOQ) &
\textcite{harris1913} \\

\ref{subsec:chinhSachCoDien} & Tính tối ưu của chính sách $(s, S)$ và các điều kiện kèm theo &
\textcite{scarf1960} \\

\ref{subsec:chinhSachCoDien} & Mô hình người bán báo (Newsvendor) một chu kỳ &
\textcite{arrow1951} \\

\ref{subsec:moHinhTonKho} & Bài toán mất đơn hàng không có nghiệm giải tích dạng đóng &
\textcite{zipkin2008} \\

\ref{subsec:phamvingoai} & Cấu trúc tồn kho đa cấp (nằm ngoài phạm vi) &
\textcite{clark1960} \\

\ref{subsec:phamvingoai} & Hiệu ứng bullwhip (nằm ngoài phạm vi) &
\textcite{lee1997} \\
\hline

\multicolumn{3}{l}{\textit{A.2. Học tăng cường}} \\
\hline
\ref{subsec:mdp} & Quá trình quyết định Markov và quy hoạch động ngẫu nhiên &
\textcite{puterman1994} \\

\ref{subsec:hamGiaTri} & Hàm giá trị, phương trình Bellman, hàm lợi thế, định lý policy gradient &
\textcite{sutton2018} \\

\ref{subsec:gae} & Ước lượng lợi thế tổng quát (GAE) và đánh đổi bias--variance &
\textcite{schulman2016gae} \\

\ref{subsec:ppo} & Vùng tin cậy trong tối ưu hóa chính sách &
\textcite{schulman2015trpo} \\

\ref{subsec:ppo} & Hàm mục tiêu cắt của PPO &
\textcite{schulman2017ppo} \\

\ref{subsec:rewardDecomposition} & Điều kiện bảo toàn chính sách tối ưu khi định hình phần thưởng &
\textcite{ng1999} \\
\hline

\multicolumn{3}{l}{\textit{A.3. Học tăng cường đa tác tử}} \\
\hline
\ref{subsec:decPomdp} & Độ phức tạp NEXP-đầy đủ của điều khiển phi tập trung &
\textcite{bernstein2002} \\

\ref{subsec:decPomdp} & Mô hình Dec-POMDP &
\textcite{oliehoek2016} \\

\ref{subsec:hocDocLap} & Học độc lập và vấn đề môi trường không dừng &
\textcite{tan1993} \\

\ref{subsec:ippo} & Thuật toán IPPO &
\textcite{dewitt2020ippo} \\

\ref{subsec:ippo} & MAPPO và khung huấn luyện tập trung, thực thi phân tán &
\textcite{yu2022mappo} \\

\ref{subsec:paramSharing} & Chia sẻ tham số giữa các tác tử đồng nhất &
\textcite{gupta2017} \\

\ref{subsec:rewardDecomposition} & Bài toán gán tín dụng đa tác tử; phần thưởng chênh lệch (so sánh với phần thưởng phân rã cục bộ) &
\textcite{wolpert2001} \\
\hline

\multicolumn{3}{l}{\textit{A.4. Ứng dụng, môi trường và giao thức đánh giá}} \\
\hline
\ref{subsec:rlTonKho} & Học tăng cường sâu cho trò chơi phân phối bia &
\textcite{oroojlooyjadid2022} \\

\ref{subsec:rlTonKho} & Đánh giá học tăng cường sâu trên các lớp bài toán tồn kho &
\textcite{gijsbrechts2022} \\

\ref{sec:moiTruong} & Chuẩn giao diện môi trường học tăng cường &
\textcite{brockman2016gym}; \textcite{towers2024gymnasium} \\

\ref{subsec:daHatGiong} & Yêu cầu đa hạt giống và tính tái lập trong đánh giá &
\textcite{henderson2018} \\

\ref{subsec:daHatGiong} & Phương pháp báo cáo kết quả có kiểm định thống kê &
\textcite{agarwal2021} \\
\hline
\end{longtable}

\vspace{0.5\baselineskip}
\noindent\textbf{Ghi chú.} Các số liệu thống kê mô tả bộ dữ liệu M5 trình bày
trong Mục~\ref{subsec:thuctrang} và Bảng~\ref{tab:phanloainhucau} do nhóm tác
giả tự tính toán từ hai tệp \texttt{sales\_train\_evaluation.csv} và
\texttt{calendar.csv}; mã nguồn tính toán được đính kèm trong kho mã của khóa
luận.
